\section{Multiple reduced incidence and its complete Hilbert boundary}
\label{sec:multiple-incidence-v161}
We remove the restriction to a single incidence point. The independence
of the incidence conditions is proved from the radicals of distinct
members of a regular pencil; it is not added as a transversality
hypothesis. The resulting modification has a global Fitting-ideal
description even though the incidence points need not be globally
labelled.

Put $M=\Sym^2V$, $P=\Pj(M)$ and $G=\Gr(2,M)$, with $n=\dim V\geq3$.
Let $Z_j\subset P$ have its symmetric determinantal scheme structure.
When $n=3$ put $Z_0=\varnothing$. Let $G_{\rm reg}$ be the open of
regular pencils and let $X_G=\Pj(\mathcal R)$ be the universal line.
The image of $X_G\times_P Z_{n-3}$ is closed, since the projection
from this incidence is proper. Remove it and call the resulting open
$G_*\subset G_{\rm reg}$. Every line over $G_*$ meets $Z_{n-2}$ in
a finite scheme: a positive-dimensional intersection would contain
the line and contradict regularity. Consequently
\[
 q_*^{\rm inc}:D_*:=X_{G_*}\times_P Z_{n-2}\longrightarrow G_*
\]
is proper and quasi-finite, hence finite. Define
\begin{equation}\label{eq:reduced-open-v161}
 U_{\rm red}=G_*\setminus
 q_*^{\rm inc}\bigl(\Supp\Omega_{D_*/G_*}\bigr),\qquad
 D=D_*\times_{G_*}U_{\rm red},\qquad q:D\to U_{\rm red}.
\end{equation}
This is open because $q_*^{\rm inc}$ is finite. Its geometric
incidences are precisely the finite reduced intersections: for a
finite algebra over $\C$, vanishing of relative differentials is
equivalent to being a product of copies of $\C$. Unlike the
single-incidence open, it allows every possible number of such points.
Set
\begin{equation}\label{eq:fitting-centre-v161}
 \mathcal Q=q_*\OO_D,\qquad
 \mathcal I=\Fitt_0(\mathcal Q),\qquad
 B_{\rm red}=\operatorname{Bl}_{\mathcal I}U_{\rm red}.
\end{equation}
The zeroth Fitting ideal of the zero module is the unit ideal.

\subsection{The general incidence calculation}
\begin{lemma}[Transverse incidence of a family of curves]
\label{lem:general-incidence-v161}
Let $S,Y$ be smooth complex varieties, let $Z\subset Y$ be smooth of
codimension $c\geq2$, and let $X\to S$ be a smooth family of embedded
curves in $Y$. Suppose $D=X\times_Y Z\to S$ is finite and unramified.
Fix $s\in S$ with incidence points $p_1,\ldots,p_k$. Write
$a_i\subset N_{Z/Y,p_i}$ for the normal image of $T_{p_i}X_s$.
Assume $a_i\ne0$ and that the joint normal evaluation map
\[
 T_sS\longrightarrow\bigoplus_{i=1}^k N_{Z/Y,p_i}/a_i
\]
is surjective. Etale locally at $s$, the incidence splits into closed
immersions $D_i\simeq\Sigma_i\hookrightarrow S$, with independent
smooth centres of codimension $c-1$. Near $p_i$ there are coordinates
in which
\begin{equation}\label{eq:general-normal-v161}
 I_D=(x_i,u_{i1},\ldots,u_{i,c-1}),\qquad
 I_{\Sigma_i}=(u_{i1},\ldots,u_{i,c-1}),
\end{equation}
where all base parameters $u_{ij}$ are independent. If
$\mathcal Q_S=q_*\OO_D$, then
\[
 \Fitt_0(\mathcal Q_S)=\prod_i I_{\Sigma_i}
                       =\bigcap_i I_{\Sigma_i}.
\]
The blow-up of this ideal is the fibre product of the individual
blow-ups. It is smooth, with etale normal-crossing exceptional boundary
and fibre $\prod_i\Pj^{c-2}$ at $s$.
\end{lemma}
\begin{proof}
Over the strict henselization at $s$, the finite algebra splits into
factors indexed by its geometric fibre. These factors and their
idempotents descend to an etale neighbourhood. For each factor the
unit map from the base ring has cokernel whose fibre at $s$ is zero:
finite base change identifies that fibre with the corresponding
one-dimensional residue algebra. Nakayama's lemma, followed by
shrinking, makes the unit map surjective. Each factor is therefore a
closed immersion. This argument uses the finite algebra, not just the
set of its geometric points.

The map of the curve tangent into the normal space is nonzero, so one
local equation of $Z$ has nonzero relative derivative. Its zero set
is etale over the base. After an etale base change choose this as a
section and use that equation as the relative coordinate $x_i$.
Modulo $x_i$, the other equations become base functions $u_{ij}$.
Their differentials give exactly the joint normal evaluation in the
statement. Surjectivity makes them independent regular parameters,
proving the normal form and smoothness of all intersections of the
centres. In particular the length-one factors in this argument are
reduced on complex geometric fibres; reducedness has not been inferred
from an arbitrary higher fibre length.

Now $\mathcal Q_S=\bigoplus_i\OO_S/I_{\Sigma_i}$. Fitting ideals of
a direct sum multiply. In the displayed independent coordinates,
intersection equals product: in a polynomial coordinate ring a
monomial in every ideal uses a variable from each disjoint group, and
the equality persists under etale flat extension. Equivalently one
may verify it in the faithfully flat completion. The simultaneous
blow-up is covered by charts in which, for each $i$, one $u_{ij}$ is
retained and the others are its multiples by new ratios. These charts
are smooth and identify the fibre product with the blow-up of the
product ideal. The latter equality also follows from the simultaneous
Rees graph lemma. There are $k$ independent exceptional parameters,
and their zero sets have normal crossings. The fibre consists of
independent choices of a line in each normal space of dimension
$c-1$, proving the last assertion.
\end{proof}

The arrangement blow-up mechanism is classical. Li's wonderful-model
Theorems~1.2 and 1.3 \cite{Li2009v161} give the normal-crossing and
product-ideal description for a building set. Here we have proved the
independent-centre case explicitly. The pencil-specific assertion
needed to apply it at every point of $U_{\rm red}$ is the following.

\begin{lemma}[Independence at distinct spectral points]
\label{lem:independent-radicals-v161}
For every $R\in U_{\rm red}$, with incidence points $p_1,\ldots,p_k$,
the joint normal evaluation of Lemma~\ref{lem:general-incidence-v161}
is surjective. In particular the local incidence centres have
codimension two and meet transversely, without any additional genericity
assumption on the other elementary divisors of $R$.
\end{lemma}
\begin{proof}
Choose a nonsingular member $Q_0$ of the pencil and write its other
members as $Q_1-\lambda Q_0$. Put $T=Q_0^{-1}Q_1$ on $V^*$.
The radicals $K_i=\ker(T-\lambda_i)$ at distinct spectral points
lie in distinct generalized eigenspaces. Their sum is direct.
Restriction therefore gives a surjection
\[
 \Sym^2V\longrightarrow\bigoplus_i\Sym^2(K_i^*):
\]
extend a basis of the direct sum of the $K_i$ to a basis of $V^*$ and
prescribe the indicated diagonal blocks of a symmetric form.
Varying $Q_1$ with $Q_0$ fixed realizes every one of these restrictions.
At a corank-two point the normal space to the rank locus is the
three-dimensional space of forms on $K_i$, with its harmless local
projective twist. The tangent of the pencil maps to the nonzero line
$a_i$, because the incidence is reduced. Quotienting by $a_i$ gives
the joint evaluation in the statement. A change of basis of the
pencil contributes only multiples of its tangent image in each
summand, so the map factors through $T_RG=\Hom(R,M/R)$. It remains
surjective after this quotient. This proves the claim also when the
tangent form on $K_i$ has rank one rather than rank two.
\end{proof}

\subsection{Identification of the full Hilbert graph}
Let $G^\circ\subset U_{\rm red}$ be the open with no corank-two
incidence. A pencil over it has a unique lifted embedded curve in
$Q=\mathrm{CQ}(V)$. Use the polarization
$H=\bigotimes_{q=1}^{n-1}H_q$, where $H_q$ is the resolved exterior
hyperplane bundle, and put $P_H(l)=\binom n2 l+1$.
Let $\Gamma_{\rm red}$ be the reduced closure of this graph in
$U_{\rm red}\times\Hilb^{P_H}(Q)$, and let
$\widehat\Gamma_{\rm red}$ be its normalization. Thus
$\widehat\Gamma_{\rm red}=\widehat G\times_G U_{\rm red}$ with the
conventions of Proposition~\ref{prop:pencil-compactification-v157}.
This equality uses localization of integral closure
\cite[Section~29.55, Lemmas~29.55.3--4]{Stacks161}; it is an open
restriction, not an arbitrary base-change assertion.

\begin{theorem}[The complete multiple-incidence modification]
\label{thm:multiple-incidence-v161}
There is a canonical congruence-equivariant isomorphism
\begin{equation}\label{eq:global-multiple-v161}
 \widehat\Gamma_{\rm red}\simeq B_{\rm red}
                =\operatorname{Bl}_{\Fitt_0(q_*\OO_D)}U_{\rm red}.
\end{equation}
The space is smooth, and its exceptional boundary has normal crossings
etale locally. Over a pencil with $k$ distinct reduced corank-two
incidences the whole scheme fibre is $(\Pj^1)^k$. A point of that fibre
specifies independently, at each $p_i$, a line through $[a_i]$ in
\[
 \pi^{-1}(p_i)=\Pj(N_{Z_{n-2}/P,p_i})\simeq\Pj^2.
\]
Its Hilbert curve is the strict transform $C_R$ of the pencil line,
with these $k$ embedded tails attached at their respective points.
It is reduced and nodal, with no embedded components. For each tail
$F_i$ and each $1\leq q<n$,
\begin{equation}\label{eq:multiple-degrees-v161}
 \deg_{F_i}H_q=\mathbf1_{q=n-1},\qquad
 \deg_{C_R}H_q=q-k\mathbf1_{q=n-1}.
\end{equation}
The incidence scheme satisfies $k\leq\lfloor n/2\rfloor$.

On a labelled etale neighbourhood the exceptional divisors $E_i$
meet like coordinate hyperplanes; every subset occurs locally, and
\[
 K_{B_{\rm red}}=\beta^*K_{U_{\rm red}}+\sum_i E_i.
\]
No global labelling, or global simplicity of each irreducible boundary
component, is asserted. The universal curve is a flat nodal family;
near its $i$th node its equation is $x_i y_i=t_i$, with $t_i=0$
the corresponding exceptional divisor. The base blow-up has its
usual effective-Cartier image-ideal universal property. These
statements classify all fibres over $U_{\rm red}$, not only the fibres
of selected smoothing arcs.
\end{theorem}
\begin{proof}
Apply the preceding two lemmas with $c=3$. Locally write
$I_i=(u_i,v_i)$. In the simultaneous blow-up chart put
$u_i=t_i$, $v_i=t_iw_i$. Near the $i$th incidence the image of the
ambient rank-locus ideal on the pulled-back universal line is
$(x_i,t_i)$. Its blow-up has charts
\[
 \OO_B[x_i,y_i]/(x_i y_i-t_i),\qquad
 \OO_B[z_i],\quad x_i=t_i z_i.
\]
The incidence points have disjoint neighbourhoods on the curve, so
these modifications do not interfere. Outside them the ideal is the
unit ideal. The local hypersurfaces are flat over $B$: they are
relative Cartier divisors in a smooth relative plane and their
nonzero defining equations remain nonzerodivisors on every fibre.
The fibrewise Cartier flatness criterion applies. These are the
standard nodal local models, so all the stated fibre properties
follow. Their total spaces are smooth in the displayed independent
coordinates.

Here is the scheme-level embedding needed for the Hilbert argument.
If $\mathcal K$ is the rank-locus ideal on the open ambient space and
$f:X_B\to P$ is evaluation, there is a surjection of graded algebras
\begin{equation}\label{eq:rees-surjection-v161}
 f^*\!\left(\bigoplus_{d\geq0}\mathcal K^d\right)
       \twoheadrightarrow
             \bigoplus_{d\geq0}(\mathcal K\OO_{X_B})^d.
\end{equation}
Its kernel removes precisely the relations annihilated on passing to
the image ideal. The target is the image Rees algebra inside
$\OO_{X_B}[T]$; it is not an assumed tensor-product identification.
Taking Proj gives a closed immersion into the base-changed ambient
blow-up, which is $Q$ over rank at least $n-2$. The universal line
is already embedded in $P\times B$. We obtain a flat closed family
of curves in $Q\times B$. This agrees with the lifted pencil on the
dense nonincident open.

In normal coordinates its exceptional map at $p_i$ is
$[x_i:t_i]\mapsto[x_i a_i+t_i b_i]$. The vectors $a_i,b_i$ are
independent, and their span determines an embedded line, whose
homogeneous ideal is its linear annihilator. Distinct directions
modulo $a_i$ therefore give distinct closed subschemes of the
exceptional plane, not merely different parametrizations of one
curve. Tails at different $p_i$ cannot be confused. The induced
proper map $B\to\Gamma_{\rm red}$ has finite geometric fibres and
is birational. It is consequently finite, and the smooth source is
the normalization of the reduced integral graph. This proves
\eqref{eq:global-multiple-v161}, including its scheme fibres. Both
maps restrict to the identity on the dense graph, giving canonicity
and equivariance and gluing the etale calculations.

The smaller minor systems have no base point, while the last minor
system has one simple base point at each $p_i$. Subtracting these
base points and adding one unit of degree on each tail gives
\eqref{eq:multiple-degrees-v161}. The determinant has vanishing order
at least two at each distinct corank-two point, so $2k\leq n$.
The nodal tree has Euler characteristic one and total $H$-degree
$\binom n2$. This proves the Hilbert polynomial used above. Finally
the product blow-up charts give independent parameters $t_i$; the
Jacobian for each codimension-two blow-up contributes one $E_i$ to
the canonical divisor. These computations also prove the local
intersection statement.
\end{proof}

\begin{corollary}[Multiplication systems at every node]
\label{cor:multi-powers-v161}
The ideal $\mathcal I$ in \eqref{eq:fitting-centre-v161} is also obtained
by taking the zeroth Fitting ideal of the pushforward of the zero
scheme of $J_{h(n-2)+1}$ on the universal line, for every $h\geq1$.
At $h=1$ this is the distinguished construction without an exponent
choice. On every tail the resolved $j$th power is the complete system
of degree
\[
 \delta_j=(j-h(n-2))_+-2(j-h(n-1))_+.
\]
The exceptional parameter space at a $k$-incidence pencil has Chow
ring $\mathbb Z[\xi_1,\ldots,\xi_k]/(\xi_1^2,\ldots,\xi_k^2)$.
In particular at a double incidence the full fibre is a smooth
quadric surface $\Pj^1\times\Pj^1$, not merely two independent
one-parameter examples.
\end{corollary}
\begin{proof}
On this rank open the universal formula gives
$J_{h(n-2)+1}=I_{n-1}$ because $I_{n-2}=\OO$. Its zero scheme is
exactly $D$. Finite pushforward and Fitting ideals give the specified
centre. The power assertion is Proposition~\ref{prop:exceptional-powers-v160}
at each disjoint tail. For clarity, a coefficient ideal generated in
degree $b$ has degree-$b$ part equal to the linear span of its actual
generators. In the residual binary ring this part is
$f^{(b-h)_+}\Sym^{b-2(b-h)_+}(\Sym^2W^*)$; division by the common
factor therefore leaves all forms, not merely a subspace with the
same base locus. The Chow ring is the projective bundle formula
applied to the identified scheme fibre $(\Pj^1)^k$.
\end{proof}
